\minitoc  % Affiche la table des matières pour ce chapitre

% ==================================================================================================================================
% Intégrale d'une fonction

\section{Intégrale d'une fonction}

\subsection{Fonction étagée et fonction intégrale}

\begin{definition}[Intégrale d'une fonction étagée]
    Soit $e$ une fonction étagée positive. On a $e = \sum_{i=1}^{n} \lambda_i 1_{A_i}$ où $A_i = e^{-1}(\{\lambda_i\})$ où $\lambda_i \in \R_+$.
    L'intégrale de $e$ par rapport à $\mu$ est de la forme :
        \[ \boxed{ \int e \; d \mu = \sum_{i=1}^{n} \lambda_i \mu (A_i) } \] 
\end{definition}

\begin{definition}[Fonction intégrale]
    Soit $f \in \overline{\mathcal{M}_+}(X)$, alors : 
        \[ \int f \; d \mu := \sup \Bigl( \Bigr\{ \int e \; d \mu : e \text{ fonction étagée positive tq } e \leq f \Bigl\}) \Bigr) \] 
    Ainsi, l'intégrale d'une fonction mesurable positive existe toujours.
    On peut donc définir la fonction intégrale telle que :
        \[ \int . \; d \mu : 
            \begin{cases}
                \overline{\mathcal{M}_+}(X) & \longrightarrow \overline{\R_+} \\ 
                f & \longmapsto \int f \; d \mu 
            \end{cases}
        \] 
\end{definition}

\begin{prop}(Propriétés de l'intégrale)
    La fonction intégrale vérifie les propriétés suivantes :
    \begin{itemize}
        \item \underline{Extension de la mesure :} $\forall A \in \mathcal{B}, \int 1 _A \; d \mu = \mu(A)$ 
        \item \underline{Linéarité positive :} $\forall f, g \in \overline{\mathcal{M}_+}(X), \forall \alpha \in \R_+$ on a :
            \[ \int (f + \alpha g) \; d \mu = \int f \; d \mu + \alpha \int g \; d \mu \]
        \item \underline{Convergence Monotone :} Pour toute suite croissante de fonctions mesurables positives (étendues) $(f_n)$, on a :
            \[ \boxed { \int \lim_{n\to\infty} f_n \; d \mu = \lim_{n\to\infty} \int f_n \; d \mu }\] 
    \end{itemize}
\end{prop}

\subsection{Intégrabilité (fonction positive)}

\begin{definition}[Fonction Intégrable (Cas Réel Positif)]
    Une fonction $f \in \overline{\mathcal{M}_+}(X)$ est dite intégrable si $\int f \;d \mu < \infty $.
\end{definition}

\begin{remark}
    Quelques remarques concernant la notation... 
    \begin{itemize}
        \item $\overline{\mathcal{L}_+^1}(X)$ est l'ensemble des fonctions $f : X \rightarrow \overline{\R_+}$ intégrables.
        \item ${\mathcal{L}_+^1}(X)$ est l'ensemble des fonctions intégrables de $X$ vers $\R_+$ mais non étendues.
        \item On remarquera que intégrable implique mesurable. La réciproques est généralement fause.
    \end{itemize}
\end{remark}

\begin{prop}[Croissance de l'intégrale]
    La fonction intégrale est croissante.
    \[ \text{i.e } \forall f,g \in \overline{\mathcal{M}_+}(X) \text{ si } f \leq g, \quad \text{alors } \int f \; d \mu \leq \int g \; d \mu \] 
\end{prop}

\begin{quote}
    \footnotesize 
    \begin{proof}
        Soient $f, g \in \overline{\mathcal{M}_+}(X)$ telles que $f \leq g$ alors $g = g-f+f$ où $g-f \geq 0 $ \\
        D'où : $ \int g \; d \mu = \int (g-f) \; d \mu + \int f \; d \mu \geq \int f \; d \mu $
    \end{proof}
    \normalsize
\end{quote}

\begin{remark}
    Pour tout $f,g \in \overline{\mathcal{M}_+}(X)$ telles que $f \leq g$ si $g$ est intégrable alors $f \in \overline{\mathcal{L}_+^1}(X)$.
\end{remark}

\subsection{Intégrabilité (fonction à valeurs complexes)}

\begin{definition}[Fonction Intégrable]
    Soit $f \in \mathcal{M}(x)$ une fonction mesurable d'un ensemble $X$ vers $\C$. 
    On dit que $f$ est intégrable si la fonction $|f| \in \mathcal{M}_+(X)$ est intégrable (ici on parle du module). \\
    On notera $\mathcal{L}^1(X)$ l'ensemble des fonction intégrables de $X$ vers $\C$.
\end{definition}

\newpage
\begin{definition}[Intégrale d'une fonction à valeurs complexes]
    Soit $f \in \mathcal{L}^1(X)$, distinguons deux cas :
    \begin{itemize}
        \item \textbf{Cas réel positif :} si $f \in \mathcal{L}^1_\R(X)$ alors soient :
            \begin{align*}
                f_+ : x \longmapsto \max (0, f(x)) \\ 
                f_+ = \sup (0, f) \\ 
                f_- : x \longmapsto \min (0, f) \\
                f_- = - \inf (0, f)
            \end{align*}
            on pose alors :
                \[ \int f \; d \mu = \int f_+ \; d \mu - \int f_- \; d \mu \] 
        \item \textbf{Cas complexe :} si $f \in \mathcal{L}^1(X)$ on pose :
            \[ \int f \; d \mu = \int \Re f \; d \mu - i \int \Im f \; d \mu \]
    \end{itemize}
\end{definition}

\begin{prop}[Propriétés de l'intégrale...2]
    Quelques propriétés sur l'intégrale...
    \begin{enumerate}
        \item \textbf{Linéarité :} La fonction intégrale sur $\mathcal{L}^1(X)$ est linéaire et $\mathcal{L}^1(X)$ est un sev de $\C$.
        \item \textbf{Croissance}
        \item \textbf{Majoration par le module :} si $f \in \mathcal{L}^1(X)$ alors $\Big| \int f \; d \mu \Big| \leq \int |f| \; d \mu$
    \end{enumerate}
\end{prop}

\begin{theorem}[Critère d'intégrabilité des fonctions bornées]
    Si $(X,\mathcal{B},\mu)$ est un espace mesuré \textbf{fini} alors toutes fonction $ f \in \mathcal{M}(X)$ \textbf{bornée} est intégrable. 
\end{theorem}

\subsection{Extension par zéro et intégrale restreinte}

\begin{definition}[Fonction extension par zéro]
    Soit $f : A \rightarrow \C$ où $A \subset X$, alors l'extension par zéro de $f$ est la fonction :
        \[ e_X(f) : 
            \begin{cases}
                X \longrightarrow \C \\
                x \longmapsto \begin{cases}
                                    f(x) \text{ si } x \in A \\
                                    0 \text{ sinon }
                                \end{cases}
            \end{cases} \] 
\end{definition}

\begin{definition}[Intégrale restreinte]
    Une fonction $f$ intégrable sur $X$ positive (resp. à valeurs complexes) est dite intégrable sur $A \subset X$ si la restriction de $f$ (resp. la restriction $|f|$) à $A$ est intégrable.
\end{definition}

\newpage 

\begin{prop}[Intégrale restreinte]
    Soient $g \in \overline{\mathcal{M}_+(A)}, A \subset X$ et $e_x$ son extension par zéro sur $X$.
    Soit $1_A$ la fonction indicatrice de $A$.
    \begin{itemize}
        \item \textbf{Intégrale et extension par zéro :}
            \[ \int_{X} e_x(g) \; d \mu = \int_{A} g \; d \mu \] 
            $g$ est donc intégrable sur A si $e_x(g)$ est intégrable sur $X$.
        \item \textbf{Intégrale restreinte et indicatrice :}
            \begin{itemize}
                \item si $f \in \mathcal{M}_+(X)$ alors $\int_{A} f \; d \mu = \int_{X} f \times 1_A \; d \mu $ 
                \item si $f \in \mathcal{M}(X)$ alors $f$ est intégrable sur $A$ si $f \times 1_A$ est intégrable sur $X$ et dans ce cas :
                    \[ \boxed{ \int_{A} f \; d \mu = \int_{X} f \times 1_A \; d \mu } \] 
            \end{itemize}
        \item \textbf{Intégrabilité et restriction :} soit $f \in \mathcal{L}^1(X)$ alors $f\restriction_A \in \mathcal{L}^1(X)$ où $A \subset X$ et 
                \[ \int_A |f| \; d \mu \leq \in_X |f| \; d \mu \]
        \item \textbf{Relation de Chasles généralisée :} Soit $(A_n)_{n\in\N}$ une suite de parties 
            mesurables deux à deux disjointes et soit $f \in \overline{\mathcal{M}_+(X)}$ alors: 
                \[ \int_{\bigcup_{n \in \N} A_n} f \; d \mu = \sum_{n=0}^{\infty} \int_{A_n} f \; d \mu \]
            Pour le cas complexe, si $f \in \overline{\mathcal{M}(X)}$  alors $f$ est intégrable sur $(A_n)_{n\in\N}$ 
            ssi $$\sum_{n=0}^{\infty} \int_{A_n} |f| \; d \mu \le \infty $$
            et dans ce cas on a l'égalité (dans $\C$) :
                \[ \int_{\bigcup_{n \in \N} A_n} f \; d \mu = \sum_{n=0}^{\infty} \int_{A_n} f \; d \mu \]
        \item \textbf{Indivisibilité des parties négligeables :} Si $N$ est une partie négligeable, toute fonction mesurable est intégrable sur $N$ et son intégrale est nulle.
                \[ \text{i.e } \int_N f \; d \mu = 0 \]
    \end{itemize}
\end{prop}
\newpage

\subsection{Intervalles non compacts et rappels}
\begin{theorem}[Intégration sur un intervalle \textbf{non compact}]
    \begin{enumerate}
        \item \textbf{Cas positif :} Soit $f \in \overline{\mathcal{M}_+([a,c[)}$ on a 
        \[ \boxed{ \int_{[a,c[} f \; d \lambda = \underset{b \in [a,c[}{\sup} \int_{[a,c[} f \; d \lambda} \]
        en particulier $ f \in \mathcal{L}^1_+([a,c[)$ :
        \begin{itemize}
            \item[\underline{ssi}] $\Bigl\{ \int_{[a,b]} f \; d \lambda : b \in [a,c[ \Bigr\}$ est majoré dans $\R_+$ 
            \item[\underline{ssi}] il existe une suite $(b_n)$ croissante de $[a,c[$ qui tend vers $c$ telle que $ (\int_{[a,bn]} f \; d \lambda)_n$ est majorée.
        \end{itemize}
        \item \textbf{Cas non positif :} Soit $ f \in \mathcal{L}^1([a,c[)$ à valeurs complexes alors :
            \[ \boxed{ \lim_{b \to c} \int_{[a,b]} d \; \lambda = \int_{[a,c[} f \; d \lambda} \] 
    \end{enumerate}
\end{theorem}

\textbf{Exemple :}

\begin{minipage}{0.4\textwidth}  % Largeur du graphique
    \begin{tikzpicture}
      \begin{axis}[
          domain=0.1:10,
          samples=200,
          width=\linewidth,  % Adapte la largeur au minipage
          height=4cm,
          xlabel=$x$,
          ylabel=$f(x)$,
          axis lines=middle,
          % xtick={0,2,4,6,8,10},
          % ytick={0,0.5,1},
          ymin=-0.5, ymax=1.5,
          % title={$f(x) = \frac{\sin(x)}{x}$ pour $x > 0$ et $f(0) = 1$},
          every axis plot/.append style={black, thick}
        ]

        % Définition des chemins pour la courbe et l'axe des abscisses
        \addplot[name path=curve, domain=0.1:10, black, thick] {sin(deg(x))/x};
        \addplot[name path=axis, domain=0.1:10] {0};

        % Zone hachurée entre la courbe et l'axe des abscisses
        \addplot [
          fill=gray, % Couleur de hachure
          opacity=0.3
        ]
        fill between [
          of=curve and axis, % Remplissage entre la courbe et l'axe des abscisses
        ];

        % Point en (0, 1) pour x = 0
        \addplot[only marks, mark=*] coordinates {(0,1)};

      \end{axis}
    \end{tikzpicture}
\end{minipage}%
\hfill
\begin{minipage}{0.55\textwidth}  % Largeur pour le texte
    $f : [a,c[ \longrightarrow \C$ peut avoir une intégrale impropre convergente sans nécessairement être intégrable au sens de Lebesque. 
    Par exemple la fonction ci-contre n'a pas d'intégrale définie au sens de Lebesgues mais au sens de Riemann oui.
    \[
    f(x) = \frac{\sin(x)}{x} \text{ lorsque } x > 0, \quad f(0) = 1.
    \]
\end{minipage}


\begin{example}[Intégrale de Riemann]
    \[ f_\alpha :
        \begin{cases}
            \R_+^* \longrightarrow \R \\
            t \longmapsto \frac{1}{t^\alpha}, \alpha \ge 0
        \end{cases}
    \] 
    \begin{enumerate}
        \item sur $[1, \infty[, f_\alpha$ est intégrable ssi $\alpha \ge 1$
        \item sur $]0, 1[, f_\alpha$ est intégrable ssi $\alpha \le 1$
    \end{enumerate}
\end{example}

\begin{example}[Intégrales de Bertrand]
    \[ f_{\alpha, \beta} :
        \begin{cases}
            \R_+^* \backslash \{1\} \longrightarrow \R \\
            t \longmapsto \frac{1}{t^\alpha | \ln t |^\beta}
        \end{cases}
    \] 
    \begin{enumerate}
        \item sur $[0, \frac{1}{2}], f_{\alpha,\beta}$ est intégrable ssi $\alpha \le 1$ ou ($\alpha = 1$ et $ \beta < 1$)
        \item sur $[2, \infty[, f_{\alpha,\beta}$ est intégrable ssi $\alpha \ge 1$ ou ($\alpha = 1$ et $ \beta > 1$)
    \end{enumerate}
\end{example}
\newpage

\subsection{Voisinnage}

\begin{definition}[Intégrabilité et voisinnage]
    Soit $f : I \longrightarrow \C$, avec $I$ un intervalle et $a \in \overline{I}$. On dit que $f$ est intégrable au voisinnage de $a$ s'il existe 
    une voisinnage $V$ de $a$ tel que :
        \[ f\restriction_{{I \cap V}} \text{ est intégrable} \] 
\end{definition}

\begin{example}
    Soit $f : [0, \infty[ \longrightarrow \C$. 
    Dire que $f$ est intégrable au voisinnage de $\infty$ signifie que $\exists c \in \R_+, f\restriction_[c, \infty[$ est intégrable.
\end{example}

\begin{definition}[Fonction localement intégrable]
    Soit $f : I \longrightarrow \C$, on dit que $f$ est localement intégrable sur $I$ si pour tout compact $K \subset I$, $f$ est intégrable sur $K$.
\end{definition}


\begin{remark}[Critère d'intégrabilité]
    Si $f : [a,b] \longrightarrow \C$ est 
    \begin{itemize}
        \item localement intégrable sur $[a,b]$
        \item intégrable au voisinnage de $b$
    \end{itemize}
    alors $f$ est intégrable sur $[a,b]$.
\end{remark}

\begin{theorem}[Changement de variable]
    Soit $\Psi : I \rightarrow I'$ de classe $\mathcal{C}^1$ et bijective. Soit $f$ une fonction positive, étendue, intégrable, alors :
        \[ \boxed{ \int_{I'} f \; d \lambda = \int_I (f \circ \Psi) \times |\Psi ' | \; d \lambda } \] 
    \[ 
		\begin{tikzcd}
			X \arrow[r, "\Psi"] \arrow[rd, "f \circ \Psi"] & I' \arrow[d, "f"]\\
                                                           & \R_+ 
		\end{tikzcd}
	\]
    Donc $f \in \mathcal{M}(X)$ est intégrable sur $I'$ ssi $(f \circ \Psi) |\Psi ' |$ est intégrable sur $I$.
\end{theorem}



\section{Théorèmes fondamentaux de l'intégrale}

\begin{theorem}[Critère d'annulation presque partout]
    Soit $f \in \overline{\mathcal{M}_+(X)}$, on a :
        \[ \int d \; \mu = 0 \Longleftrightarrow f = 0 \text{ p.p} \]
\end{theorem}

\begin{theorem}[Critère de finitude presque partout]
    Soit $f \in \overline{\mathcal{L}_+^1(X)}$, alors $f \le \infty $ p.p 
\end{theorem}

\newpage

\begin{theorem}[Inversion série intégrale(Cas positif)]
    Soit $(f_n)_{n\in\N}$ une suite de fonction dans $\overline{\mathcal{M}_+(X)}$ 
        \[ \text{alors } \int \sum_{n \geq 0} f_n \; d \mu = \sum_{n \geq 0} \int f_n \; d \mu \] 
    \textbf{Attention :} Cette propriété n'est pas à confondre avec la linéarité car ici, nous sommes en présence d'une limite de somme de fonctions.
\end{theorem}

\begin{theorem}[Convergence Dominée]
    Soit $(f_n)_{n\in\N}$ une suite de fonctions dans $\overline{\mathcal{M}(X)}$.
    Supposons que $(f_n)_{n\in\N}$ \underline{converge simplement} vers $f \in \mathcal{M}(X)$ 
    \[ \boxed{
    \begin{aligned}
        & \text{\underline{si}} \quad  \exists g \in \mathcal{L}_+^1(X) \text{ telle que } \forall n \in \N, |f_n| \leq g \\
        & \text{\underline{alors}} \quad \forall n \in \N, f_n \text{ est intégrable et }  f \text{ est intégrable sur} X
    \end{aligned} }
    \]
    \[ \text{et } \int \Big| f_n - f \Big| \; d \mu \underset{n \to \infty}{\longrightarrow} 0 \quad \Longleftrightarrow \quad  \int f_n \; d \mu \underset{n \to \infty}{\longrightarrow} \int f \; d \mu \] 
    \[ \text{i.e } \lim_{n \to \infty} \int f_n \; d \mu = \int \lim_{n \to \infty} f_n \; d \mu \] 
    ... dans $\C$
\end{theorem}

\begin{theorem}[Inversion série-intégrale (Cas complexe)]
    Soit $ \sum f_n$ une série de fonctions telle que $ \forall n \in \N, f_n \in \mathcal{L}^1(X)$ converge absoluments pour tout $x \in X$
    $ \Big( \text{ie} \quad \forall x \in X, \sum_{n \in \N} \Big| f_n (x) \Big| \le \infty \Big)$
    \[ \boxed{
    \begin{aligned}
        & \text{\underline{si}} \quad  \sum_{n \in \N} \int \Big| f_n \Big| \; d \mu \le \infty \\
        & \text{\underline{alors}} \quad \int \sum_{n \in \N} f_n \; d \mu = \sum_{n \in \N} \int f_n \; d \mu 
    \end{aligned} }
    \]
    et $\sum_{n \in \N} f_n$ est intégrable sur $X$ (en tant que fonction à valeurs complexes).
\end{theorem}

\begin{quote}
\footnotesize
\begin{proof}
    Posons $g_n = \sum_{k=0}^{n} f_k$. Comme $\sum f_n$ converge, $(g_n)$ converge simplement vers $\sum_{n=0}^{\infty} f_n = f $

    Soit $n \in \N$, on sait que $ |g_n| \leq \Big| \sum_{k=0}^{n} f_k \Big| \leq \sum_{k=0}^{n} | f_k | \leq \sum_{n=0}^{\infty} f_n := g $
    Donc $g$ est intégrable. On peut appliquer le théorème de convergence dominée. \\
    \underline{Alors} $ \sum_{n=0}^{\infty} f_n$ est intégrable et :
    \begin{align*}
        \lim_{n \to \infty} \int g_n \; d \mu &= \int \sum_{n=0}^{\infty} f_n \; d \mu \\
        \lim_{n\to\infty} \int \sum_{k=0}^{n} f_k \; d \mu &= \int \sum_{n=0}^{\infty} f_n \; d \mu \\
        \lim_{n\to\infty} \sum_{k=0}^{n} \int f_k \; d \mu &= \int \sum_{n=0}^{\infty} f_n \; d \mu \\
        \sum_{n=0}^{\infty} \int f_n \; d \mu &= \int \sum_{n=0}^{\infty} f_n \; d \mu 
    \end{align*}
\end{proof}
\normalsize
\end{quote}
